%===========================================================
% TCC — FGV/EMAp
% Predição e Recomendação de Escalações no Cartola FC
% Gabriel Rodrigues de Deus Abad
% Orientador: Prof. Moacyr Alvim Horta Barbosa da Silva
%===========================================================

\documentclass[12pt, a4paper, oneside]{article}

%------------------------------------------------------------
% Pacotes
%------------------------------------------------------------
\usepackage[utf8]{inputenc}
\usepackage[T1]{fontenc}
\usepackage[brazil]{babel}
\usepackage{lmodern}
\usepackage{microtype}
\usepackage{geometry}
\usepackage{setspace}
\usepackage{indentfirst}
\usepackage{amsmath, amssymb}
\usepackage{mathtools}
\usepackage{booktabs}
\usepackage{array}
\usepackage{multirow}
\usepackage{graphicx}
\usepackage{float}
\usepackage{subcaption}
\usepackage{xcolor}
\usepackage{hyperref}
\usepackage{cleveref}
\usepackage{natbib}
\usepackage{titlesec}
\usepackage{fancyhdr}
\usepackage{caption}
\usepackage{enumitem}
\usepackage{tcolorbox}

%------------------------------------------------------------
% Layout
%------------------------------------------------------------
\geometry{a4paper, top=2.5cm, bottom=2.5cm, left=3.0cm, right=2.5cm}
\setstretch{1.3}
\setlength{\parindent}{1.25cm}
\setlength{\headheight}{15pt}

%------------------------------------------------------------
% Cores
%------------------------------------------------------------
\definecolor{azulfgv}{RGB}{0,51,102}
\definecolor{verdedest}{RGB}{26,122,74}
\definecolor{cinzaclaro}{RGB}{245,247,250}

%------------------------------------------------------------
% Seções
%------------------------------------------------------------
\titleformat{\section}
  {\large\bfseries\color{azulfgv}}{\thesection.}{0.6em}{}
  [\vspace{-0.3em}\textcolor{azulfgv}{\rule{\linewidth}{0.8pt}}]
\titleformat{\subsection}
  {\normalsize\bfseries\color{azulfgv}}{\thesubsection.}{0.5em}{}
\titleformat{\subsubsection}
  {\normalsize\bfseries}{\thesubsubsection.}{0.5em}{}

%------------------------------------------------------------
% Cabeçalho/rodapé
%------------------------------------------------------------
\pagestyle{fancy}
\fancyhf{}
\fancyhead[L]{\small\textcolor{azulfgv}{\textbf{FGV/EMAp}}}
\fancyhead[R]{\small Gabriel Rodrigues de Deus Abad}
\fancyfoot[C]{\small\thepage}
\renewcommand{\headrulewidth}{0.6pt}

%------------------------------------------------------------
% Caixa de implicação
%------------------------------------------------------------
\tcbuselibrary{skins,breakable}
\newtcolorbox{insightbox}{
  enhanced, colback=cinzaclaro, colframe=azulfgv,
  fonttitle=\bfseries\small, title=Implicação para Modelagem,
  left=4pt, right=4pt, top=4pt, bottom=4pt,
  boxrule=0.6pt, breakable
}

%------------------------------------------------------------
% Hyperref
%------------------------------------------------------------
\hypersetup{
  colorlinks=true, linkcolor=azulfgv,
  citecolor=verdedest, urlcolor=azulfgv,
  pdftitle={Predicao e Recomendacao de Escalacoes no Cartola FC},
  pdfauthor={Gabriel Rodrigues de Deus Abad}
}

%============================================================
\begin{document}
%============================================================

%------------------------------------------------------------
% CAPA
%------------------------------------------------------------
\thispagestyle{empty}
\begin{center}
  {\large\bfseries\color{azulfgv} FUNDAÇÃO GETULIO VARGAS}\\[0.3em]
  {\normalsize\color{azulfgv} ESCOLA DE MATEMÁTICA APLICADA -- FGV/EMAp}\\[0.3em]
  {\normalsize\color{azulfgv} CURSO DE GRADUAÇÃO EM MATEMÁTICA APLICADA}\\[2.5cm]
  {\Large\bfseries\color{azulfgv}
    Predição e Recomendação de Escalações no Cartola FC}\\[0.6em]
  {\large\itshape
    Combinando Aprendizado de Máquina e Ciência de Redes\\
    para Maximizar a Pontuação por Rodada}\\[2.5cm]
  {\normalsize
    \begin{tabular}{rl}
      \textbf{Autor:}      & Gabriel Rodrigues de Deus Abad \\[0.3em]
      \textbf{Orientador:} & Prof.\ Moacyr Alvim Horta Barbosa da Silva \\[0.3em]
      \textbf{Curso:}      & Graduação em Matemática Aplicada \\[0.3em]
      \textbf{Data:}       & 2026 \\
    \end{tabular}}
\end{center}
\vfill\newpage

%------------------------------------------------------------
% RESUMO
%------------------------------------------------------------
\thispagestyle{empty}
\begin{center}{\large\bfseries Resumo}\end{center}
\vspace{0.5em}

O Cartola FC é o maior \textit{fantasy game} de futebol do mundo em número de usuários
ativos. A escolha de quais jogadores escalar a cada rodada constitui um problema de
otimização combinatória com incerteza: a pontuação futura de cada atleta depende de
desempenho recente, qualidade do adversário, mando de campo e disponibilidade para jogo.
Este trabalho propõe um sistema \textit{end-to-end} que, a partir dos dados históricos
das 18 primeiras rodadas do Campeonato Brasileiro 2026, prevê a pontuação individual de
cada atleta e recomenda a escalação de maior pontuação esperada dentro das restrições do
jogo. A metodologia combina: (i)~engenharia de \textit{features} temporais com janelas
deslizantes sobre os scouts individuais; (ii)~métricas de ciência de redes:
PageRank ponderado e sistema Elo, calculadas sobre um grafo cumulativo de resultados,
para quantificar a força relativa dos adversários; (iii)~comparação de modelos preditivos
(Ridge, XGBoost, LightGBM e LSTM) sob protocolo de validação temporal
\textit{walk-forward}; e (iv)~resolução do problema de escalação via Programação Linear
Inteira (PLI). A análise exploratória revela que o status do atleta é a variável de
filtragem mais crítica: atletas ``Prováveis'' apresentam taxa de participação de 86,4\%
e pontuação média 6,3$\times$ superior à de atletas em ``Dúvida''. As \textit{features}
de desempenho histórico de curto prazo dominam o poder preditivo pré-treino, com
correlação de Spearman com o \textit{target} acima de 0,45.

\vspace{1em}\noindent
\textbf{Palavras-chave:} Fantasy Sports; Cartola FC; Aprendizado de Máquina;
Ciência de Redes; PageRank; Programação Linear Inteira; Séries Temporais.

\newpage

%------------------------------------------------------------
% SUMÁRIO
%------------------------------------------------------------
\thispagestyle{empty}
\tableofcontents
\newpage

%============================================================
\setcounter{page}{1}

% ============================================================
\section{Introdução}
\label{sec:intro}
% ============================================================

O Cartola FC é um jogo de \textit{fantasy} esportivo em que cada usuário monta
semanalmente um time de doze integrantes (onze jogadores de campo e um técnico)
respeitando restrições de posição e orçamento em cartoletas. A pontuação de cada atleta
é calculada com base em eventos reais da partida: gols (8,0~pts), assistências (5,0~pts),
defesas exclusivas de goleiros (1,3~pt), desarmes (1,5~pt) e penalidades como cartões
amarelos ($-$1,0~pt), entre outros. O objetivo de cada participante é maximizar a
pontuação total do time por rodada, tornando a decisão de escalação um problema de
otimização combinatória com incerteza.

Atualmente, a maioria dos usuários apoia suas escolhas em análises subjetivas ou
tendências de redes sociais. A disponibilidade de dados granulares por rodada via API
pública do Cartola FC cria uma oportunidade concreta para a aplicação de métodos
quantitativos rigorosos.

\subsection{Objetivos}

O objetivo geral é desenvolver um sistema preditivo e de recomendação capaz de sugerir,
para cada rodada, a escalação de maior pontuação esperada dentro das restrições do
Cartola FC. Os objetivos específicos são:

\begin{enumerate}[leftmargin=*, label=(\roman*)]
  \item Realizar análise exploratória dos dados das 18 primeiras rodadas do Brasileirão 2026;
  \item Construir um conjunto de \textit{features} temporais, contextuais e estruturais
        sem vazamento de dados;
  \item Integrar métricas de ciência de redes, PageRank ponderado e Elo, como
        \textit{features} de força do adversário;
  \item Comparar modelos preditivos de pontuação individual com protocolo
        \textit{walk-forward};
  \item Formular e resolver o problema de escalação via Programação Linear Inteira.
\end{enumerate}

\subsection{Contribuições}

A contribuição metodológica central é a integração de métricas de redes complexas,
construídas sobre um grafo cumulativo de resultados, como \textit{features} em modelos
preditivos de pontuação individual no Cartola FC. O trabalho documenta um
\textit{pipeline end-to-end} reproduzível, desde a coleta via API até a recomendação
final, com protocolo de validação temporal rigoroso que elimina qualquer vazamento de
dados futuros.

% ============================================================
\section{Revisão Bibliográfica}
\label{sec:revisao}
% ============================================================

A literatura de predição em \textit{fantasy sports} é dominada pelo contexto
norte-americano (NFL). \citet{beal2020} demonstraram que modelos de \textit{ensemble}
em particular \textit{gradient boosting}, que superam médias históricas quando
\textit{features} de contexto (adversário, mando, fadiga) são incorporadas, resultado
que orienta diretamente a engenharia de \textit{features} proposta neste trabalho.
No futebol europeu, \citet{pappalardo2019} construíram um conjunto de dados de eventos
de alta resolução e mostraram que contribuições individuais vão além dos scouts
tradicionais.

Para ranqueamento via redes, \citet{page1999} introduziram o PageRank, cuja extensão
ao contexto esportivo é natural: uma equipe é forte se venceu equipes fortes.
O sistema de ratings Elo, proposto por \citet{elo1978} para o xadrez e adaptado ao
futebol, oferece alternativa dinâmica amplamente adotada em plataformas como
FiveThirtyEight. Na etapa de otimização, o problema de escalação se assemelha ao
\textit{Multiple-choice Knapsack Problem} \citep{kellerer2004}; a formulação via PLI
garante solução ótima para as dimensões do problema. Os modelos preditivos incluem
o XGBoost \citep{chen2016} como referência estado-da-arte em dados tabulares e redes
LSTM \citep{hochreiter1997} para captura explícita de dependências temporais.

% ============================================================
\section{Dados}
\label{sec:dados}
% ============================================================

\subsection{Fontes e Coleta}

Os dados provêm de duas fontes referentes às 18 primeiras rodadas do Brasileirão 2026:
(i)~\textbf{API pública do Cartola FC}: por rodada, dados de cada atleta escalável:
pontuação real (\texttt{pontos\_num}), preço (\texttt{preco\_num}), 18 scouts individuais
e status de disponibilidade (\texttt{status\_id}); (ii)~\textbf{Resultados das partidas}
(\texttt{games.csv}): placar oficial com identificação de mandante e visitante.

\subsection{Caracterização}

\begin{table}[H]
\centering
\caption{Características do dataset bruto}
\label{tab:dataset}
\begin{tabular}{lcc}
\toprule
\textbf{Dimensão} & \textbf{players.csv} & \textbf{games.csv}\\
\midrule
Registros totais  & 13.215 & 179 \\
Rodadas           & 18     & 18  \\
Atletas únicos    & 845    & {-} \\
Clubes            & 20     & 20  \\
Variáveis         & 28     &  7  \\
Valores ausentes  & 0      &  0  \\
\bottomrule
\end{tabular}
\end{table}

\subsection{Amostra de Treino}
\label{subsec:amostra}

A modelagem restringe-se à \textbf{amostra de treino}: atletas com status ``Provável''
(status\_id\,=\,7) que efetivamente entraram em campo (\texttt{entrou\_em\_campo}~=~1).
Esta restrição é motivada empiricamente na Seção~\ref{subsec:status}: a diferença de
pontuação média entre ``Provável'' e ``Dúvida'' é de 6,3$\times$
(3,68 vs.\ 0,58~pts), e a taxa de participação cai de 86,4\% para 24,9\%.
Dos 13.215 registros totais, 3.294 compõem a amostra de treino.

% ============================================================
\section{Metodologia}
\label{sec:metodologia}
% ============================================================

O \textit{pipeline} segue cinco etapas sequenciais:

\medskip
\noindent
\textbf{Coleta \& Pré-proc.}
$\to$ \textbf{EDA}
$\to$ \textbf{Eng.\ de \textit{Features}}
$\to$ \textbf{Modelagem}
$\to$ \textbf{Otimização}
\medskip

\subsection{Engenharia de \textit{Features}}
\label{subsec:features}

O dataset final \texttt{df\_features} tem uma linha por par (atleta,~rodada) e
$\approx$65 colunas. O protocolo \textit{walk-forward} garante que toda \textit{feature}
da rodada $t$ usa exclusivamente dados das rodadas $1,\ldots,t-1$, eliminando
\textit{data leakage}. As \textit{features} são organizadas em cinco blocos:

\subsubsection*{Bloco 1: Desempenho Histórico do Atleta}
Médias móveis de pontuação com janelas $k \in \{3,5,10\}$ rodadas jogadas
(\texttt{pts\_avg\_3r}/\texttt{5r}/\texttt{10r}), desvio padrão móvel
(\texttt{pts\_std\_3r}/\texttt{5r}), máximo e mediana nas últimas 5~rodadas, e
tendência linear (\texttt{pts\_trend\_5r}), coeficiente angular de regressão
de pontuação sobre rodada nas últimas 5~aparições. Para cada scout relevante
(G, A, FT, FF, FD, FS, FC, CA, DS, DE, SG, GS), médias móveis de 3 e 5~rodadas.
Completam o bloco: proporção de rodadas com pontuação positiva (\texttt{prop\_jogou})
e sequência consecutiva de rodadas acima da própria média histórica
(\texttt{pts\_acima\_media\_streak}).

\subsubsection*{Bloco 2: Contexto da Partida}
Mando de campo binário (\texttt{is\_mandante}); força do adversário via PageRank
e Elo calculados sobre o grafo acumulado até $t-1$
(\texttt{adv\_pagerank}, \texttt{adv\_elo}); médias de gols marcados e sofridos
pelo adversário nas últimas 5~rodadas; histórico \textit{head-to-head} do atleta
contra aquele adversário (\texttt{h2h\_pts\_avg}); e fase do campeonato
(\texttt{fase\_campeonato}).

\subsubsection*{Bloco 3: Preço e Mercado}
Preço atual (\texttt{preco\_num}), variação acumulada nas últimas 3~rodadas
(\texttt{preco\_variacao\_3r}) e índice custo-benefício histórico
(\texttt{cb\_avg\_5r}).

\subsubsection*{Bloco 4: \textit{Features} do Time}
Forma nas últimas 5~rodadas em pontos de tabela (\texttt{time\_forma\_5r}),
médias de gols marcados/sofridos pelo clube, e PageRank/Elo do próprio time.

\subsubsection*{Bloco 5: Posição e Status}
Codificação \textit{one-hot} de posição; \texttt{status\_id}; e probabilidade
empírica de participação dado o status (\texttt{prob\_jogar}), calculada sobre
o histórico disponível até $t-1$.

\subsection{Ciência de Redes}
\label{subsec:redes}

Modelamos o campeonato como grafo direcionado e ponderado
$G_t = (V, E_t)$: 20~nós (clubes), arestas do perdedor ao vencedor ponderadas pelo
saldo de gols, e empates como arestas bidirecionais com peso~0,5. O grafo é
reconstruído cumulativamente a cada rodada.

\paragraph{PageRank.}
Aplicamos o algoritmo de \citet{page1999} com fator de amortecimento $d = 0{,}85$:
\begin{equation}
  \mathrm{PR}(u) = \frac{1-d}{N} + d \sum_{v \to u}
  \frac{\mathrm{PR}(v)\,w(v,u)}{\displaystyle\sum_{x} w(v,x)},
  \label{eq:pagerank}
\end{equation}
onde $N = 20$ e $w(v,u)$ é o peso da aresta. O valor $\mathrm{PR}(\text{adversário})$
codifica de forma não-linear e transitiva a dificuldade do confronto.

\paragraph{Elo.}
Rating dinâmico inicializado em 1500 para todos os clubes e atualizado após cada
partida:
\begin{equation}
  R'_A = R_A + K\bigl(S_A - \mathbb{E}[S_A]\bigr),
  \qquad
  \mathbb{E}[S_A] = \frac{1}{1 + 10^{(R_B - R_A)/400}},
  \label{eq:elo}
\end{equation}
onde $K$ é o fator de sensibilidade (ajustado empiricamente), $S_A \in \{0;\,0{,}5;\,1\}$
é o resultado real e $\mathbb{E}[S_A]$ é o resultado esperado. O Elo do clube
\emph{antes} da rodada $t$ entra como \textit{feature}.

\subsection{Modelos Preditivos}
\label{subsec:modelos}

Formulamos o problema como regressão supervisionada: dado o vetor
$\mathbf{x}_{i,t}$, estimar $\hat{y}_{i,t} \approx y_{i,t}$ (pontuação real).
O protocolo \textit{walk-forward} treina o modelo nas rodadas $4,\ldots,t-1$
e avalia na rodada $t$, repetindo para $t = 5,\ldots,18$.

Os modelos comparados são:
\begin{enumerate}[leftmargin=*, label=(\roman*)]
  \item \textbf{Baseline}: média ponderada das últimas 5~rodadas do atleta;
  \item \textbf{Ridge / Lasso}: regressão linear com regularização $L_2$ / $L_1$,
        $\alpha$ ajustado por validação cruzada temporal;
  \item \textbf{Random Forest}: 300~árvores, \textit{feature importance} via
        impureza de Gini;
  \item \textbf{XGBoost / LightGBM}: \textit{gradient boosting} com
        hiperparâmetros otimizados via Optuna (50 tentativas, otimização bayesiana);
  \item \textbf{LSTM}: 2~camadas de 64~unidades, \textit{dropout}~0,2,
        janelas deslizantes de 5~rodadas por atleta.
\end{enumerate}

Os modelos são treinados \textbf{separadamente por posição} (motivação detalhada
na Seção~\ref{subsec:feat_pos}).
Métricas de avaliação: RMSE, MAE, $R^2$ e correlação de Spearman entre pontuação
prevista e real (esta diretamente relevante para a qualidade do ranking de escalação).

\subsection{Otimização de Escalação}
\label{subsec:otimizacao}

Dadas as pontuações previstas $\{\hat{y}_{i,t}\}$, o problema de escalação é:

\begin{align}
  \max_{\mathbf{x}} \quad
  & \sum_{i} \hat{y}_{i,t}\,x_i \label{eq:obj}\\
  \text{s.a.} \quad
  & \sum_{i} c_i\,x_i \leq C \tag{orçamento}\\
  & \sum_{i:\,p_i=k} x_i = n_k,\quad
    k \in \{\text{Gol, Lat, Zag, Mei, Ata, Téc}\} \tag{posições}\\
  & \sum_{i:\,\text{clube}(i)=j} x_i \leq 5,\quad \forall j \tag{concentração}\\
  & x_i = 0,\quad \forall i : \text{status}(i) \neq \text{Provável} \tag{disponib.}
\end{align}

com $n_k$ vagas por posição (1 Goleiro, 2 Laterais, 2 Zagueiros, 3 Meias,
3 Atacantes, 1 Técnico) e $x_i \in \{0,1\}$. Resolvido via solver CBC / PuLP
em tempo sub-segundo para $\approx$700 atletas disponíveis por rodada.

% ============================================================
\section{Resultados: EDA e Engenharia de \textit{Features}}
\label{sec:resultados}
% ============================================================

\subsection{Distribuição de Pontuações por Posição}

\begin{figure}[H]
  \centering
  \includegraphics[width=\linewidth]{fig1_dist_posicao.png}
  \caption{Distribuição de pontuações por posição, restrita à amostra de treino
           (Provável e em campo). (A)~Violinplot de densidade. (B)~Boxplot com
           entalhes representando o IC~95\% da mediana.}
  \label{fig:dist_pos}
\end{figure}

\begin{table}[H]
\centering
\caption{Estatísticas descritivas por posição (amostra de treino)}
\label{tab:stats_pos}
\begin{tabular}{lccccc}
\toprule
\textbf{Posição} & \textbf{N} & \textbf{Média} & \textbf{Mediana}
                 & \textbf{Desvio P.} & \textbf{Máx.}\\
\midrule
Goleiro   & 311  & 4,08 & 3,40 & 3,95 & 15,90\\
Lateral   & 475  & 4,79 & 4,30 & 4,04 & 19,00\\
Zagueiro  & 555  & 3,32 & 2,40 & 3,47 & 17,40\\
Meia      & 950  & 3,78 & 2,70 & 3,84 & 19,90\\
Atacante  & 691  & 4,92 & 3,20 & 4,87 & 26,00\\
Técnico   & 312  & 5,22 & 4,55 & 2,35 & 12,76\\
\bottomrule
\end{tabular}
\end{table}

Atacantes exibem o maior desvio padrão (4,87~pts), confirmando a maior
imprevisibilidade desta posição. Técnicos têm a maior média (5,22~pts) com
variância reduzida, pois sua pontuação depende do resultado coletivo do time.

\subsection{Análise Temporal e Autocorrelação}

\begin{figure}[H]
  \centering
  \includegraphics[width=\linewidth]{fig2_temporal.png}
  \caption{Evolução temporal da pontuação média por posição. (A)~Séries temporais
           por rodada. (B)~Heatmap posição $\times$ rodada.}
  \label{fig:temporal}
\end{figure}

O coeficiente de variação (CV) permanece entre 86\% e 108\% em todas as 18~rodadas,
refletindo a natureza assimétrica dos scouts. A análise de ACF/PACF das séries
individuais (calculada apenas sobre rodadas Provável e em campo, sem imputação de
zeros) indica dependência temporal predominantemente de lag~1.

\begin{figure}[H]
  \centering
  \includegraphics[width=0.92\linewidth]{fig3_acf_pacf.png}
  \caption{ACF e PACF para séries individuais de três atletas representativos:
           Pedro (Atacante), Danilo (Meia) e Lucas Arcanjo (Goleiro).
           Bandas: IC~95\% sob hipótese de ruído branco.}
  \label{fig:acf}
\end{figure}

\begin{insightbox}
Ausência de autocorrelação significativa em lags superiores a 2 justifica janelas
de médias móveis de $k \in \{3,5\}$ rodadas como mecanismo principal de captura
de tendência. Janelas longas (10~rodadas) são mantidas como \textit{feature}
complementar, mas com menor peso esperado.
\end{insightbox}

\subsection{Status e Disponibilidade}
\label{subsec:status}

\begin{figure}[H]
  \centering
  \includegraphics[width=\linewidth]{fig7_status_disponibilidade.png}
  \caption{Análise de status: (A)~volume por categoria; (B)~pontuação média;
           (C)~taxa de participação em campo (\texttt{entrou\_em\_campo});
           (D)~distribuição de pontuação por status.}
  \label{fig:status}
\end{figure}

\begin{table}[H]
\centering
\caption{Pontuação e participação por status (base completa)}
\label{tab:status}
\begin{tabular}{lcccc}
\toprule
\textbf{Status} & \textbf{N} & \textbf{Média (pts)}
                & \textbf{\% Zeros} & \textbf{\% Particip.}\\
\midrule
Provável (7)        & 3.812 & 3,68 & 16,4\% & 86,4\%\\
Dúvida (6)          & 7.809 & 0,58 & 79,1\% & 24,9\%\\
Suspenso (5)        &   936 & 0,10 & 96,8\% &  4,1\%\\
Lesionado (3)       &   197 & 1,99 &  6,6\% & 93,4\%\\
Não Participará (2) &   461 & 2,58 & 31,9\% & 74,2\%\\
\bottomrule
\end{tabular}
\end{table}

A diferença de pontuação entre ``Provável'' e ``Dúvida'' é de 6,3$\times$,
tornando o status a variável de filtragem mais crítica: o modelo treina exclusivamente
na amostra Provável, e o otimizador exclui atletas com status diferente.

\subsection{Mando de Campo}

\begin{figure}[H]
  \centering
  \includegraphics[width=\linewidth]{fig6_mandante_visitante.png}
  \caption{Efeito do mando de campo. (A)~Boxplot global ($t = 7{,}07$, $p < 0{,}001$).
           (B)~Diferença por posição (Mandante~$-$~Visitante).
           (C)~Violinplot por posição e mando.}
  \label{fig:mando}
\end{figure}

O efeito mandante é significativo em 4 de 5 posições de linha ($\alpha = 0{,}05$),
com Laterais apresentando o maior ganho absoluto (+1,19~pts, $p = 0{,}001$). Goleiros
são a exceção ($p = 0{,}17$, n.s.), possivelmente porque o desempenho defensivo
coletivo dilui o efeito do mando nesta posição.

\begin{table}[H]
\centering
\caption{Efeito do mando de campo por posição (amostra de treino)}
\label{tab:mando}
\begin{tabular}{lcccc}
\toprule
\textbf{Posição} & $\bar{y}_{\text{mand.}}$ & $\bar{y}_{\text{visit.}}$
                 & $\Delta$ & \textbf{Signif.}\\
\midrule
Goleiro   & 4,39 & 3,77 & $+$0,62 & n.s.\\
Lateral   & 5,38 & 4,20 & $+$1,19 & **\\
Zagueiro  & 3,84 & 2,79 & $+$1,05 & ***\\
Meia      & 4,14 & 3,39 & $+$0,75 & **\\
Atacante  & 5,45 & 4,37 & $+$1,08 & **\\
\bottomrule
\multicolumn{5}{l}{\small *** $p<0{,}001$, ** $p<0{,}01$, n.s.\ não significativo (teste $t$ de Welch)}
\end{tabular}
\end{table}

\subsection{Análise do \texttt{df\_features}}

\subsubsection{Poder Preditivo Pré-Treino}

\begin{figure}[H]
  \centering
  \includegraphics[width=0.78\linewidth]{fig2_spearman_target.png}
  \caption{Correlação de Spearman entre cada \textit{feature} e
           \texttt{pontos\_num} (rodada~$\geq 4$). Barras vermelhas indicam
           correlação negativa. Ordenado por $|\rho|$ decrescente.}
  \label{fig:spearman}
\end{figure}

As \textit{features} de desempenho histórico dominam: \texttt{pts\_avg\_10r},
\texttt{pts\_avg\_5r} e \texttt{pts\_max\_5r} atingem $\rho \approx 0{,}47$
(Figura~\ref{fig:spearman}). As métricas de ciência de redes
(\texttt{adv\_pagerank}, \texttt{adv\_elo}) apresentam $\rho \approx 0$ na análise
univariada; o efeito da força do adversário é inerentemente interativo (um atacante
forte pontua mais contra defesas fracas) e deve ser capturado pelos modelos
não-lineares na etapa de modelagem.

\subsubsection{Multicolinearidade}

\begin{figure}[H]
  \centering
  \includegraphics[width=\linewidth]{fig3_multicolinearidade.png}
  \caption{Heatmap de correlação de Spearman entre \textit{features}
           (A) e zoom no bloco de pontuação e preço (B).}
  \label{fig:multicol}
\end{figure}

Há alta correlação entre médias móveis de diferentes janelas:
$\rho(\texttt{pts\_avg\_3r},\,\texttt{pts\_avg\_5r}) = 0{,}91$ e
$\rho(\texttt{pts\_avg\_5r},\,\texttt{pts\_avg\_10r}) = 0{,}93$.

\begin{insightbox}
Multicolinearidade elevada ($\rho > 0{,}84$) entre janelas de pontuação motiva
regularização forte nos modelos lineares (Ridge/Lasso) e favorece modelos de árvore,
que são naturalmente robustos por selecionar apenas uma \textit{feature} redundante
por nó de decisão. Esta é uma das justificativas para esperar que XGBoost e LightGBM
superem a regressão Ridge.
\end{insightbox}

\subsubsection{Distribuição das \textit{Features} por Posição}
\label{subsec:feat_pos}

\begin{figure}[H]
  \centering
  \includegraphics[width=\linewidth]{fig7_features_por_posicao.png}
  \caption{Distribuição de \textit{features} de scout por posição (rodada~$\geq 4$).
           \texttt{DE\_avg\_5r} é concentrada em Goleiros; \texttt{G\_avg\_5r}
           é dominante em Atacantes.}
  \label{fig:feat_pos}
\end{figure}

As distribuições diferem radicalmente entre posições:
\texttt{DE\_avg\_5r} (defesas) é quase zero em todas as posições exceto Goleiro;
\texttt{GS\_avg\_5r} (gols sofridos) é exclusivo desta posição.
Esta heterogeneidade estrutural justifica o \textbf{treinamento de modelos separados
por posição}: um modelo único não captura simultaneamente os padrões de um Goleiro
(dominado por defesas e \textit{clean sheets}) e de um Atacante (dominado por gols
e finalizações).

\subsubsection{Estabilidade das Métricas de Redes}

\begin{figure}[H]
  \centering
  \includegraphics[width=\linewidth]{fig5_pagerank_elo.png}
  \caption{Evolução do PageRank (A) e do Elo (B) ao longo das 18~rodadas.
           Rodada~1 não possui PageRank (grafo vazio); todos os clubes
           iniciam com Elo~=~1500.}
  \label{fig:pr_elo}
\end{figure}

O PageRank exibe alta instabilidade nas primeiras 4--5~rodadas (grafo esparso),
convergindo progressivamente. O Elo separa com clareza os clubes dominantes:
Palmeiras e Flamengo acima de 1.580 a partir da rodada~10, e os times em queda.
Esta estabilidade diferencial sugere que o Elo é mais informativo que o PageRank
nas rodadas iniciais.

% ============================================================
\section{Experimentos e Resultados de Modelagem}
\label{sec:experimentos}
% ============================================================

\textit{Esta seção será preenchida após a implementação dos modelos. Os experimentos
seguirão o protocolo walk-forward (rodadas~5--18) descrito na
Seção~\ref{subsec:modelos}.}

Os resultados a reportar incluem:
\begin{itemize}[leftmargin=*]
  \item Tabela comparativa de RMSE, MAE, $R^2$ e Spearman por modelo e por posição;
  \item \textit{Feature importance} do XGBoost por posição;
  \item Pontuação das escalações sugeridas vs.\ escalação ótima (retrospecto) e
        média dos usuários da plataforma;
  \item Percentil de desempenho das escalações na distribuição de todos os
        participantes do Cartola FC.
\end{itemize}

% ============================================================
\section{Conclusão}
\label{sec:conclusao}
% ============================================================

\textit{Esta seção será completada após a finalização dos experimentos.}

Esta fase do trabalho estabeleceu as bases analíticas e a estrutura completa do
\textit{pipeline}. Os principais achados são:

\begin{enumerate}[leftmargin=*, label=(\roman*)]
  \item A restrição à amostra Provável é empiricamente justificada: diferença de
        6,3$\times$ na pontuação e queda de 86,4\% para 24,9\% na taxa de
        participação ao comparar ``Provável'' com ``Dúvida'';
  \item \textit{Features} de curto prazo dominam o poder preditivo pré-treino
        ($\rho_\text{Spearman} \approx 0{,}47$), consistente com dependência
        de lag~1 identificada na ACF;
  \item Alta multicolinearidade entre janelas ($\rho > 0{,}84$) favorece modelos
        de \textit{gradient boosting} sobre regressão linear;
  \item Heterogeneidade posicional nas distribuições de scouts justifica modelos
        separados por posição;
  \item Métricas de redes (PageRank, Elo) têm correlação univariada próxima de zero,
        mas contribuição interativa é esperada nos modelos não-lineares.
\end{enumerate}

%============================================================
% REFERÊNCIAS
%============================================================
\newpage
\bibliographystyle{plainnat}
\bibliography{referencias}

\end{document}
